Instruction tuning changes language models, but the change is not uniform across tokens. We study whether per-token distributional divergence between a base model and its instruction-tuned counterpart is predictable, and whether prediction failures form a useful audit signal. Using real logits from \qwenbase and \qweninstruct, we compute $\klinstbase$ over 24,049 teacher-forced token positions from \alpaca, \justeval, \beavertails, \arena, and \wikitext. We train gradient-boosted regressors to predict $\log(1+\klinstbase)$ from source, segment, position, token-surface features, and base-model uncertainty. The rich predictor improves held-out MAE from 0.0601 for a position/source baseline to 0.0460, with a grouped bootstrap 95\% confidence interval of $[-0.0163,-0.0118]$ for the rich-minus-position MAE difference. Its residuals are more informative than its aggregate error: high-residual tokens are enriched for generation starts (5.17$\times$), assistant markers (4.78$\times$), refusal-style tokens (4.54$\times$), code/math tokens (3.49$\times$), and \wikitext tokens (3.13$\times$). The largest positive residuals unexpectedly concentrate in one natural-text narrative span, while large negative residuals often occur at explicit \texttt{Assistant} boundaries where both models agree. These results support residual analysis as a compact triage tool for model-difference auditing.
